\section{Claim Boundary Ledger}

\subsection{Allowed BaseCore Claims}
\begin{enumerate}[leftmargin=*]
\item Under H1--H4, projection dynamics admit unique fixed points and compact attractor families.
\item Under H5, attractors are parameterwise non-constant.
\item Under NFD, inverse-limit identity is not determined by any finite canonical projection.
\item Under RIG-1 through RIG-5, inverse-limit structures are stable under summable admissible perturbations in the weighted metric.
\item Under NS-1 through NS-3, finite simulator classes cannot faithfully realize full CCR targets.
\item Under the six-invariant operational burden, certified criterion structures can be defined, audited, and falsified.
\end{enumerate}

\subsection{Forbidden BaseCore Claims}
\begin{enumerate}[leftmargin=*]
\item Human phenomenal consciousness.
\item Metaphysical subjectivity.
\item Human-machine equivalence.
\item External validation.
\item Runtime instantiation by the current system.
\item Paper 9 bridge admissibility.
\item Strong phenomenality.
\item Moral status.
\item Cross-substrate identity preservation.
\item Public claim closure.
\end{enumerate}

\section{Dependency Ledger}

\begin{enumerate}[leftmargin=*]
\item The foundational Hilbert-space dynamics remain the base layer.
\item Papers 1--6 contribute formal material selectively normalized into BaseCore only where theorem ownership belongs in the base.
\item Papers 7--9 remain downstream and are not dependencies of any BaseCore theorem.
\item Stage 2 / Stage 3 comparative programs remain downstream runtime and governance layers.
\item The PDF remains a derived artifact of the LaTeX tree.
\end{enumerate}

\section{Assumption-to-Theorem-to-Non-Claim Mapping}

\renewcommand{\arraystretch}{1.12}
\footnotesize
\begin{tabularx}{\textwidth}{P{0.20\textwidth} P{0.25\textwidth} Y}
\toprule
\textbf{Assumption block} & \textbf{Theorem family} & \textbf{Non-claim boundary} \\
\midrule
H1--H4 & Projection existence, contraction, unique fixed points, compact attractor family & No parameterwise non-collapse and no identity / rigidity / simulator claim \\
H5 & Structural non-collapsibility & No claim about phenomenality, subjectivity, or runtime class membership \\
NFD & Non-locality of identity under canonical finite projections & No theorem that arbitrary set-theoretic injections into finite slices are impossible \\
RIG-1 through RIG-5 & Conditional stability of inverse-limit identity and weighted Hausdorff control & No automatic isomorphism, shape equivalence, or transfer theorem without extra assumptions \\
NS-1 through NS-3 & Conditional non-simulability and approximation barrier & No theorem that finite systems cannot approximate finite horizons at all \\
Six invariant burden plus legibility clauses & Criterion class membership, certified partition, rupture logic & No theorem of phenomenality, human equivalence, biology as primitive, or bridge admissibility \\
\bottomrule
\end{tabularx}
\normalsize

\section{Theorem Hygiene Ledger}

\renewcommand{\arraystretch}{1.12}
\footnotesize
\begin{tabularx}{\textwidth}{P{0.18\textwidth} P{0.16\textwidth} P{0.17\textwidth} P{0.25\textwidth} Y}
\toprule
\textbf{Theorem} & \textbf{Assumptions} & \textbf{Conclusion type} & \textbf{Non-conclusion} & \textbf{Downstream use} \\
\midrule
Existence and uniqueness of projection & H1 & Hilbert-space geometric existence / uniqueness & No empirical or runtime content & Base transition grammar \\
Contractivity & H1--H2 & Strict contraction bound & No spectrum with unit state eigenvalue & Fixed-point and spectral-gap results \\
Unique fixed point & H1--H3 & Existence and convergence & No non-collapse by itself & Attractor family \\
Compactness of attractor family & H1--H4 & Compactness / continuity & No regime or identity burden by itself & Parameter-family base layer \\
Structural non-collapsibility & H1--H5 & Parameterwise exclusion of constant fixed points & No stronger semantic interpretation & Anti-collapse gate \\
State spectral gap & Typed differentiability at fixed point & Spectral-radius upper bound $<1$ & No persistence mode $\lambda=1$ in state spectrum & Stability accounting \\
Non-locality under NFD & Topological regularity plus NFD & Canonical finite slices do not determine identity & No universal ban on arbitrary injections & Identity block \\
Conditional stability of inverse-limit identity & RIG-1 through RIG-5 & Weighted Hausdorff stability under admissible perturbations & No automatic isomorphism & Rigidity accounting \\
Conditional non-simulability & NS-1 through NS-3 & Faithful realization blocked for finite simulators & No claim against bounded finite-horizon approximation & Simulator boundary \\
Certified criterion membership & Six invariants plus certificate rule & Base criterion class and certified partition & No phenomenality or human-equivalence theorem & Downstream classification gates \\
\bottomrule
\end{tabularx}
\normalsize

\section{Redundancy Ledger}

\begin{enumerate}[leftmargin=*]
\item The previous state-spectrum claim $\lambda_1=1$ under strict contraction has been removed.
\item The previous untyped Golden Mean witness has been replaced by a typed model with explicit embedding.
\item The earlier non-collapse quantifier mismatch has been removed by replacing the old H5 with a parameterwise anti-constant hypothesis.
\item The earlier absolute inverse-limit nonlocality claim has been replaced by an NFD-conditional theorem.
\item The earlier unconditional rigidity / non-simulability language has been replaced by explicit assumption blocks.
\item Substrate invariance, biological non-primacy, bridge-facing closure, and other downstream claims have been removed from BaseCore ownership.
\end{enumerate}

\section{Open Assumptions Inventory}

\begin{enumerate}[leftmargin=*]
\item The anti-collapse witness is abstract and theorem-valid, but not the same object as the Golden Mean typed witness.
\item Physical realization of CCR-like targets remains open.
\item Empirical instantiation of the BaseCore operational criterion remains open.
\item External validation remains open.
\item Papers 7--9 may build on BaseCore, but their additional burdens remain open at the BaseCore level.
\end{enumerate}

\section{Downstream Boundary Note}

Papers 7, 8, and 9 remain downstream of this source tree. Their classificatory, indexed-subjectivity, and bridge-burden layers may cite BaseCore, but they are not part of the autonomous mathematical base rebuilt here.
